% !TEX root = ../../../main.tex

\subsection{词级分类模型结构与训练策略}

在获得 X.npy、y.npy 和 labels.json 后，训练脚本加载全部 $20\times232$ 特征样本。
其中，X.npy 保存词级视频样本的时序特征矩阵，y.npy 保存对应的类别编号，labels.json 保存类别编号与 Gloss 标签之间的映射关系。
系统根据 labels.json 中的标签数量确定模型输出层的分类类别数。
为了保证训练集、验证集和测试集中各类别分布相对均衡，系统采用按类别分层划分策略对样本进行划分。
这种方式可以避免部分类别只出现在训练集或测试集中，从而提高模型评估结果的可靠性。

\WhuAutoFigure{0.82\textwidth}{0.42\textheight}{figures/chapter06/word-classifier-structure.png}{词级分类模型结构图}{fig:chapter06-word-classifier-structure}

如图~\ref{fig:chapter06-word-classifier-structure} 所示，词级分类模型采用轻量级双向 GRU 时序分类结构。
模型输入为 $20\times232$ 的 plus 特征序列，其中 20 表示重采样后的帧数，232 表示每一帧的特征维度。
输入序列首先经过 Masking 层，用于处理可能存在的零值填充帧。
随后，模型使用两层双向 GRU 提取时间序列中的动作变化信息。
第一层双向 GRU 设置 96 个隐单元，并返回完整时间序列，以保留每个时间步的中间表示。
第二层双向 GRU 设置 64 个隐单元，并输出定长向量，用于后续分类判断。
每层 GRU 后接 Layer Normalization，以稳定中间特征分布并改善训练过程。

双向 GRU 结构能够同时利用动作序列的前向信息和后向信息。
对于词级手语动作而言，某些类别可能在动作起始阶段较为相似，但在动作结束姿态或运动方向上存在差异。
因此，双向时序建模有助于模型综合考虑完整动作片段中的前后文变化。
与直接使用单帧特征分类相比，该结构能够更好地学习手型变化、手部运动轨迹和上肢姿态变化之间的时序关系。
同时，相较于更复杂的视频端到端模型，基于关键点特征的双向 GRU 模型参数规模较小，更适合当前小词表实验和系统原型验证。

在完成 GRU 时序特征提取后，模型通过全连接层进一步压缩特征并完成类别判别。
第一层全连接层设置 128 个神经元，并使用 ReLU 激活函数增强非线性表达能力。
随后加入 Dropout 层，以降低小样本训练场景下的过拟合风险。
第二层全连接层设置 64 个神经元，同样使用 ReLU 激活函数和 Dropout 正则化。
最后，输出层的神经元数量等于当前小词表类别数，并通过 Softmax 函数输出各 Gloss 类别的概率分布。
模型最终将概率最高的类别作为 Top-1 预测结果，同时保留 Top-3 候选结果用于后续句子识别与语义修正。

\begin{longtable}{L{0.24\textwidth} L{0.24\textwidth} L{0.40\textwidth}}
  \caption{词级分类模型训练参数表}
  \label{tab:chapter06-word-model-train-params}\\
  \toprule
  参数 & 取值 & 说明 \\
  \midrule
  输入形状 & $20\times232$ & 单个词级视频样本的时序特征 \\
  模型结构 & BiGRU + Dense & 基于关键点序列的时序分类模型 \\
  优化器 & Adam & 用于模型参数优化 \\
  学习率 & $1\times10^{-3}$ & 初始学习率 \\
  损失函数 & sparse categorical crossentropy & 多分类任务损失函数 \\
  batch\_size & 8 & 适配当前小样本训练规模 \\
  epochs & 120 & 最大训练轮数 \\
  dropout & 0.35 & 降低模型过拟合风险 \\
  评价指标 & accuracy / \path{top3_accuracy} & 评估 Top-1 与 Top-3 分类效果 \\
  \bottomrule
\end{longtable}

如表~\ref{tab:chapter06-word-model-train-params} 所示，模型训练使用 Adam 优化器和稀疏多分类交叉熵损失函数。
由于标签数组 y.npy 中保存的是整数类别编号，因此损失函数采用 sparse categorical crossentropy，而不需要将标签额外转换为 one-hot 编码。
训练过程中，accuracy 用于衡量模型 Top-1 预测是否正确，\path{top3_accuracy} 用于衡量真实类别是否出现在概率最高的前三个候选类别中。
对于后续句子视频识别任务而言，Top-3 候选结果具有重要意义。
即使某个词级片段的 Top-1 预测不完全准确，只要真实词汇仍然位于候选集合中，后续的语义修正模块仍有机会结合上下文对结果进行调整。

从整体训练策略看，词级分类模型并不直接处理原始视频像素，而是基于前一阶段构建的 $20\times232$ plus 特征进行训练。
这种方式降低了背景、光照和服装等视觉因素对模型的干扰，使模型更加关注人体姿态、手部结构和动作变化本身。
同时，固定长度输入、分层数据划分、Dropout 正则化和 Top-3 评价指标共同构成了当前句子视频识别实验的训练策略。
该模型的输出结果不仅用于单个词级视频的分类评估，也为后续连续视频中的滑动窗口预测、候选 Gloss 生成和语义修正提供基础。